\section{Instrumento de Coleta de Dados e Compilação das Entrevistas de Imersão}
\label{apendice:questionario}

Apresenta-se a seguir a estrutura do roteiro semiestruturado aplicado presencialmente na fase de imersão do \emph{Design Thinking}, seguido da compilação e tabulação das respostas obtidas junto a dez cidadãos abordados em nós intermodais e praças da cidade de São Paulo.

% -------------------------------------------------------------
% PARTE I: ESTRUTURA DO FORMULÁRIO / QUESTIONÁRIO
% -------------------------------------------------------------
\subsection{Roteiro do Questionário de Imersão}
\begin{enumerate}[leftmargin=*]


    \item \textbf{[Pergunta-Filtro Eliminatória] Você esteve ou está ativamente buscando imóvel para alugar ou comprar na cidade de São Paulo nos últimos 12 meses?}
    \begin{itemize}[leftmargin=1.5em]
        \item[] (\quad) \textbf{Sim} $\rightarrow$ \emph{Prosseguir para a Questão 2.}
        \item[] (\quad) \textbf{Não} $\rightarrow$ \emph{Entrevista encerrada com agradecimento.}
    \end{itemize}

    \item \textbf{Perfil Geral e Escopo Habitacional:}
    \begin{itemize}[leftmargin=1.5em]
        \item Faixa etária e ocupação principal.
        \item Tipo de transição habitacional pretendida (aluguel individual, locação familiar, república estudantil ou compra de imóvel próprio).
        \item Regiões ou bairros de interesse preferencial na capital.
    \end{itemize}

    \item \textbf{Critérios Decisórios e Segurança Pública:}
    \begin{itemize}[leftmargin=1.5em]
        \item Além do valor financeiro do aluguel ou condomínio, quais fatores urbanos são mais decisivos na escolha do endereço?
        \item O que mais gera sensação de segurança no seu trajeto diário a pé: a presença de muros altos e portarias privadas ou ruas com fluxo constante de pedestres, comércio ativo e boa iluminação pública?
        \item Você já desistiu de visitar ou alugar um imóvel devido à ``fama'' de perigoso de um determinado bairro? De onde veio essa percepção?
    \end{itemize}

    \item \textbf{Transparência Pública e Acesso à Informação:}
    \begin{itemize}[leftmargin=1.5em]
        \item Você sabia que a Prefeitura Municipal de São Paulo mantém o portal GeoSampa, disponibilizando mapas abertos de equipamentos públicos (saúde, parques, zoneamento)?
        \item Você sabia que a Secretaria de Segurança Pública de São Paulo (SSP-SP) publica tabelas mensais com todas as ocorrências policiais registradas por Distrito Policial? Já tentou consultar essas bases ao escolher um imóvel?
    \end{itemize}

    \item \textbf{Avaliação da Proposta de Tecnologia Cívica (GeoHabita):}
    \begin{itemize}[leftmargin=1.5em]
        \item Em uma escala de 1 (nada útil) a 5 (essencial), qual o grau de utilidade de uma ferramenta pública e gratuita que informe o panorama de segurança e a distância a pé até UBSs, metrô e praças a partir da simples digitação do CEP?
        \item Você considera relevante que a ferramenta apresente dados em tabelas textuais com recursos de acessibilidade (modo de alto contraste, leitor de tela para pessoas com baixa visão)?
    \end{itemize}

    \item \textbf{Consentimento e Identificação do Participante:}
    \begin{itemize}[leftmargin=1.5em]
        \item[] (\quad) Desejo me identificar formalmente no relatório da pesquisa (Nome/E-mail).
        \item[] (\quad) Prefiro participar de forma \textbf{estritamente anônima}, ciente de que nenhum dado pessoal será exposto e que as respostas serão associadas exclusivamente a um código alfanumérico identificador (ex: Munícipe 1). Autorizo a guarda confidencial do meu contato exclusivamente para convocação das sessões de teste e feedback do protótipo.
    \end{itemize}
\end{enumerate}

\vspace{1em}

% -------------------------------------------------------------
% PARTE II: COMPILAÇÃO DAS 10 RESPOSTAS COLETADAS
% -------------------------------------------------------------
\subsection{Compilação Qualitativa das Respostas Coletadas}

Todas as pessoas selecionadas que responderam afirmativamente à pergunta-filtro (Questão 1) e optaram expressamente pela preservação do anonimato (Questão 6), sendo identificadas ordinalmente de Munícipe 1 a Munícipe 10.

\begin{table}[!htbp]
\centering
\small
\renewcommand{\arraystretch}{1.25}
\begin{tabular}{|c|p{3.2cm}|p{4.2cm}|c|c|c|}
\hline
\textbf{ID} & \textbf{Perfil do Entrevistado} & \textbf{Critérios Principais} & \textbf{GeoSampa?} & \textbf{SSP-SP?} & \textbf{Nota} \\
\hline
M1 & Estudante (21 anos), busca quarto próx. metrô & Transporte noturno, iluminação e comércio no caminho & Não & Não & 5 \\
\hline
M2 & Autônoma (35 anos), aluguel com 1 filho & Proximidade de creche/UBS, calçadas movimentadas & Não & Não & 5 \\
\hline
M3 & Dev. de software (28 anos), compra 1º imóvel & Acesso a eixos de mobilidade, praças para pets & Não & Não & 4 \\
\hline
M4 & Aposentada (64 anos), aluguel & Acessibilidade física, UBS próxima, iluminação & Não & Não & 3 \\
\hline
M5 & Atendente (24 anos), aluguel compartilhado & Pontos de ônibus seguros, fluxo de pedestres & Não & Não & 5 \\
\hline
M6 & Professor estadual (39 anos), aluguel & Corredor de ônibus, preço realista, escolas & Não & Não & 4 \\
\hline
M7 & Casal com bebê (31/29 anos), aluguel 2 dorm. & Parques públicos, calçadas largas, tranquilidade & Não & Não & 5 \\
\hline
M8 & Pós-graduando (26 anos, vindo do interior) & Segurança nos trajetos a pé, proximidade à USP & Não & Não & 5 \\
\hline
M9 & Comerciário (44 anos), compra imóvel popular & Linhas de trem/metrô, custo benefício do bairro & Não & Não & 4 \\
\hline
M10 & Designer (30 anos), usuária de bike/transporte & Ciclovias, vitalidade urbana, vida cultural & Não & Não & 5 \\
\hline
\end{tabular}
\caption{Quadro sintético dos entrevistados na fase de imersão empírica.}
\label{tab:resumo_entrevistados}
\end{table}

\vspace{0.5em}

\noindent\textbf{Síntese dos Relatos Coletados por Participante:}

\begin{description}[leftmargin=1.2em, itemsep=0.8em]
    \item[Munícipe 1 (Abordado no Terminal Pinheiros):] Jovem universitário buscando locação no Butantã ou Pinheiros. Relatou que sua maior aflição é o retorno das aulas após as 22h. Declarou preferir ruas com padarias e transeuntes a condomínios fechados em vias ermas: \emph{``O prédio pode ser blindado, mas o trajeto da estação até o portão é onde o assalto acontece''}. Desconhecia GeoSampa e os relatórios da SSP. Atribuiu nota 5 ao projeto.

    \item[Munícipe 2 (Abordada no Terminal Barra Funda):] Mãe solo buscando aluguel na Zona Norte/Oeste. Prioriza a proximidade com serviços de saúde e educação infantil. Declarou ter receio de bairros periféricos por causa do noticiário da televisão, mas admite não ter dados reais: \emph{``A imobiliária só fala maravilhas, a TV só fala de tragédia; a gente fica no escuro''}. Nunca ouviu falar de dados abertos do governo. Atribuiu nota 5.

    \item[Munícipe 3 (Abordado no Largo da Batata):] Profissional de tecnologia avaliando compra de apartamento na Zona Leste/Centro Expandido. Reclamou da opacidade dos anúncios de corretoras, que vendem segurança com base em câmeras e portarias remotas. Não conhecia os portais públicos municipais. Destacou a importância de haver métricas de áreas verdes no entorno. Atribuiu nota 4.

    \item[Munícipe 4 (Abordada no entorno da Praça da República):] Idosa que procura locação próxima a equipamentos públicos de saúde na região central. Declarou que não sabia da existência do GeoSampa e aprovou a inclusão de recursos de acessibilidade e leitor de tela em tabelas textuais. Atribuiu nota 3.

    \item[Munícipe 5 (Abordada no Terminal Barra Funda):] Trabalhadora do setor de serviços buscando moradia no eixo da Linha 3-Vermelha. Apontou que a volta para casa em pontos de ônibus escuros é seu principal fator de estresse diário. Atribuiu nota 5.

    \item[Munícipe 6 (Abordado no Terminal Pinheiros):] Professor da rede básica procurando imóvel na Zona Sul. Afirmou que o custo do aluguel o empurra para regiões que ele não conhece bem e que sofre com a falta de referências confiáveis sobre os bairros: \emph{``A gente fica refém de comentários vagos em grupos de redes sociais''}. Nunca consultou a SSP-SP por achar que os dados eram de uso restrito da polícia. Atribuiu nota 4.

    \item[Munícipe 7 (Abordados na Praça Roosevelt):] Casal à procura de imóvel de dois dormitórios para alugar na Bela Vista ou Cambuci. Enfatizaram a busca por calçadas caminháveis e praças arborizadas para passear com o carrinho do filho. Desconheciam qualquer canal de dados abertos da PMSP. Atribuíram nota 5.

    \item[Munícipe 8 (Abordado nas imediações do Butantã):] Estudante recém-chegado de outro estado buscando aluguel compartilhado. Declarou completa desorientação sobre quais distritos são tranquilos para pedestres: \emph{``Quem vem de fora não sabe nem o que é DP. Se eu puder colocar só o CEP do anúncio e entender o bairro, resolve metade das minhas dúvidas''}. Não conhecia as bases da SSP-SP ou GeoSampa. Atribuiu nota 5.

    \item[Munícipe 9 (Abordado na Estação Santo Amaro):] Comerciário buscando financiamento pelo programa habitacional na Zona Sul. Relatou que a segurança de transporte coletivo é prioritária sobre sistemas privados de alarme. Afirmou que as informações do governo parecem feitas para técnicos e burocratas, não para trabalhadores. Atribuiu nota 4.

    \item[Munícipe 10 (Abordada na Praça Benedito Calixto):] Moradora em busca de locação em bairros mistos (comércio e moradia integrados). Enfatizou que locais com bares e restaurantes na calçada transmitem muito mais segurança do que quarteirões com muros extensos e guaritas. Destacou a relevância de cruzar dados de mobilidade cicloviária e saúde com estatísticas de furto/roubo. Atribuiu nota 5.
\end{description}